\subsection{Certified temporal non-integrability on the induced quotient}

The high-frequency argument is carried out on the stable quotient of a
finite-connector Young magnet.  This point is essential: the billiard map is
invertible, whereas the quotient return map is a uniformly expanding map with
countably many inverse branches.  The spectral data of the quotient and of
the anisotropic collision operator agree on the leading eigenspace.

Let
\[
 c_0=(0,0),\qquad c_1=(1,0),\qquad
 c_2=(1/2,\sqrt3/2),
\]
and let \(c=(1/2,\sqrt3/6)\) be the center of the triangular gap.  The three
inward normal collision points are
\[
 p_0=(\sqrt3R/2,R/2),\quad
 p_1=(1-\sqrt3R/2,R/2),\quad
 p_2=(1/2,\sqrt3/2-R).
\]
They form an equilateral triangle of side
\[
 d_R=1-\sqrt3R\in[1-47\sqrt3/100,1-9\sqrt3/20],
\]
which is uniformly positive.  The two oriented words
\[
 w_+=(0,1,2,0),\qquad w_-=(0,2,1,0)
\]
are regular specular return polygons.  At every internal vertex the two
adjacent unit directions have the inward disk normal as their angle bisector.

For endpoints with arclength coordinates \((r,r')\) on a common small arc
about \(p_0\), let \(\mathscr L_{R,\pm}(r,r')\) be the total Euclidean length
of the \(w_\pm\) polygon after the two internal collision coordinates have
been eliminated by the specular critical equations.

\begin{lemma}[Uniform polygon generating functions]
\label{lem:r3-a2-polygon}
There are a radius-independent endpoint interval \(J\), a neighborhood of the
symmetric collision polygon, and constants \(c_H,C_H>0\) such that the
internal critical points defining \(\mathscr L_{R,\pm}\) exist uniquely and
are \(C^3\) in \((R,r,r')\).  The internal Hessian of the polygonal length
satisfies
\[
 c_H I\le D^2_{\rm int}\mathscr L_{R,\pm}\le C_H I.
\]
At the symmetric diagonal endpoint,
\[
 \partial_{r'}\mathscr L_{R,+}(0,0)
 -\partial_{r'}\mathscr L_{R,-}(0,0)=1.
\]
Consequently, after shrinking \(J\) once,
\[
 \inf_{R\in\mathcal I}\inf_{(r,r')\in J^2}
 \left|\partial_{r'}(\mathscr L_{R,+}-
 \mathscr L_{R,-})(r,r')\right|\ge\frac34.
\]
\end{lemma}

\begin{proof}
For a segment joining two strictly convex boundary components, the mixed
second derivative of its length generating function is
\[
 -\frac{\cos\varphi\cos\varphi'}{\tau},
\]
and each diagonal second derivative is the sum of the adjacent free-flight
term and the curvature term \(R^{-1}\cos^{-1}\varphi\).  On the two triangular
polygons all flight lengths lie in a compact subinterval of \((0,\infty)\)
and all collision angles stay a fixed distance from grazing, uniformly for
\(R\in\mathcal I\).  Strict dispersing curvature therefore makes the two by
two internal Hessian uniformly positive.  The implicit-function theorem gives
the stated critical points and regularity.

After the internal variables are eliminated, the envelope theorem says that
the terminal derivative is the tangential component of the incoming unit
momentum.  At \(p_0\) the positively oriented last leg has direction
\((-1/2,-\sqrt3/2)\), the negatively oriented last leg has direction
\((-1,0)\), and the positive unit tangent to the disk is
\((-1/2,\sqrt3/2)\).  Their tangential components are respectively
\(-1/2\) and \(1/2\); reversing the endpoint convention changes both signs
and leaves the absolute difference equal to one.  Uniform continuity on the
compact radius interval gives the \(3/4\) bound.
\end{proof}

Choose one common regular connector word \(c_*\) from the terminal arc back to
the base of the magnet.  Precompose both triangular words with \(m\) copies of
\(c_*\), where \(m\) is fixed below.  Let \(h_{R,+},h_{R,-}\) be the resulting
inverse branches of the stable-quotient return map and let
\(\mathcal R_{R,+},\mathcal R_{R,-}\) be their return roofs.  Stable holonomy
identifies their terminal endpoints.

\begin{lemma}[Uniform temporal-distance derivative]
\label{lem:r3-a2-uni}
There are a fixed connector length \(m\), a common quotient interval \(J_0\),
and \(c_{\rm UNI}>0\) such that the temporal distance
\[
 \Psi_R(x)=\mathcal R_{R,+}(h_{R,+}x)
 -\mathcal R_{R,-}(h_{R,-}x)
 +\sum_{n\ge1}
 \left(\tau_R(F_R^{-n}h_{R,+}x)-
       \tau_R(F_R^{-n}h_{R,-}x)\right)
\]
is \(C^1\) on \(J_0\) and satisfies
\[
 \inf_{R\in\mathcal I}\inf_{x\in J_0}|\Psi_R'(x)|
 \ge c_{\rm UNI}.
\]
One may take \(c_{\rm UNI}=1/4\).
\end{lemma}

\begin{proof}
The first difference is the reduced polygonal generating-function difference
of Lemma~\ref{lem:r3-a2-polygon}, transported by stable holonomy.  The
holonomy Jacobian differs from one by a uniformly bounded factor and can be
made arbitrarily close to one by shrinking the common interval.  The common
connector makes the derivative of the remaining past tail at most
\[
 C\Lambda_*^{-m}\sum_{n\ge0}\Lambda_*^{-n}.
\]
Choose \(m\) so this is at most \(1/4\), and then shrink the interval so the
holonomy error is at most \(1/4\).  The \(3/4\) polygon derivative then leaves
a lower bound \(1/4\).  All choices are uniform because the geometry and
hyperbolicity constants are uniform in \(R\).
\end{proof}

\begin{theorem}[Joint aperiodicity and high-frequency contraction]
\label{thm:r3-a2-high}
If \((u,t)\in[-\pi,\pi]^2\times\mathbb R\) and
\(\mathcal L_{R,iu,it}\) has peripheral spectral radius one, then
\(u=0\) and \(t=0\).  For \(|t|\ge1\),
\[
 \|\mathcal L_{R,iu,it}^n\|_{\mathcal B_R\to\mathcal B_{R,w}}
 \le C|t|^A\exp\left(-\frac{cn}{\log(2+|t|)}\right)
\]
with constants uniform in \(R\) and \(u\).
\end{theorem}

\begin{proof}
A peripheral mode descends to a unit-modulus eigenfunction on the induced
stable quotient.  The periodic-orbit identity makes
\(u\cdot S_n\kappa_R-tS_n\tau_R\) constant modulo \(2\pi\) on every full
branch.  Differentiating the two branch identities in
Lemma~\ref{lem:r3-a2-uni} forces \(t=0\).  With \(t=0\), the open branches
having displacements \(\pm e_1,\pm e_2\) force
\(u\in2\pi\mathbb Z^2\), hence \(u=0\) in the fundamental domain.

For the quantitative estimate, partition the quotient unstable interval into
cells of length comparable with \(|t|^{-1}\).  On every return to \(J_0\),
the two branch phases have derivative separation at least
\(c_{\rm UNI}|t|\); van der Corput cancellation removes a fixed fraction of
the paired standard-family mass.  The growth lemma returns a fixed fraction
of every family to \(J_0\) in \(O(\log(2+|t|))\) iterates.  Iteration gives
the displayed contraction.  The anisotropic quotient/intertwining maps carry
it back to the collision operator.  Compact nonzero frequency annuli are
handled by upper semicontinuity and the already-proved absence of peripheral
modes.
\end{proof}
